% Part: history
% Chapter: biographies
% Section: julia-robinson

\documentclass[../../../include/open-logic-section]{subfiles}

\begin{document}

\olfileid{his}{bio}{rob}

\olsection{جوليا روبنسون}

\olphoto{robinson-julia}{Julia Robinson}

جوليا بومان روبنسون رياضية أمريكية، جلّ شهرتها من بحوثها في مسائل القرار،
ومن أعظم آثارها نصيبها في حل مسألة هيلبرت العاشرة. وُلدت في سانت لويس من
ولاية ميزوري يوم 8 ديسمبر 1919، وكانت، فيما تروي عن صباها، مولعة بالأعداد
منذ نعومة أظفارها \citep[4]{Reid1986}. ولما بلغت التاسعة أصابتها الحمى
القرمزية، ثم تتابعت عليها نوبات الحمى الروماتيزمية، فلبثت زمناً طويلاً
طريحة الفراش وتأخر درسها. وقد استدركت بمعونة معلّمين خاصين ما فاتها من
التعليم، إلا أن المرض خلّف في بدنها أثراً لازمها بقية عمرها.

ولم تمنعها شدائد طفولتها أن تفرغ من المدرسة الثانوية حائزة جوائز شتى في
الرياضيات والعلوم. وابتدأت تعليمها الجامعي بكلية ولاية سان دييغو، ثم
انتقلت في عامها الأخير إلى جامعة كاليفورنيا ببركلي، وهناك لقيت الرياضي
رافائيل روبنسون وتأثرت به؛ فتوثقت الصداقة بينهما ثم تزوجا سنة 1941. ولما
كانت زوجة رجل من هيئة التدريس مُنعت أن تدرّس في قسم الرياضيات ببركلي.
وظلت تحضر مقررات الرياضيات مستمعة، وإن كانت ترجُو أن تدع الجامعة وتنشئ
أسرة. ثم أصابها التهاب رئوي بعيد زواجها، وأخبرها الأطباء أن الحمى
الروماتيزمية التي عانتها في الصغر قد أورثت قلبها نسيجاً ندبياً كثيراً؛
حتى ظن الطبيب أنها لا تجاوز الأربعين، وأشار عليها ألا تلد
\citep[13]{Reid1986}.

ومكثت روبنسون بعد ذلك حيناً طويلاً في كآبة، ثم عزمت أخيراً على استئناف
الرياضيات، فعادت إلى بركلي وأتمت الدكتوراه سنة 1948 بإشراف ألفرد تارسكي.
وكان تارسكي قد برهن أن نظرية الرتبة الأولى للأعداد الحقيقية قابلة للقرار،
وكان يلزم من عمل غودل أن نظيرتها في الأعداد الطبيعية غير قابلة له؛ أما
قابلية نظرية الرتبة الأولى للأعداد النسبية للقرار فبقيت من كبار المسائل
المفتوحة. فأثبتت روبنسون في أطروحتها \citeyearpar{Robinson1949} أن تلك
النظرية غير قابلة للقرار.

ثم غلب عليها الاشتغال بمسائل القرار، فقصدت مسألة هيلبرت العاشرة، وهي واحدة
من ثلاث وعشرين مسألة رياضية شهيرة عرضها ديفيد هيلبرت سنة 1900. وفحوى
العاشرة: هل توجد خوارزمية تقرر في زمن متناه هل لمعادلة كثيرة حدود ذات معاملات
صحيحة، كالمعادلة $3x^2 - 2y +3 = 0$، حل في الأعداد الصحيحة أم لا؟ وتسمى
هذه \emph{المسائل الديوفانتية}. وبعد نجاحها الأول تعاونت روبنسون مع مارتن
ديفيس وهيلاري بوتنام، وكانا مشتغلين بالمسألة نفسها؛ فأبان الثلاثة أن
المسائل الديوفانتية الأسية، التي يجوز أن تقع فيها المجاهيل في مواضع
الأسس، غير قابلة للقرار. وأبانوا أيضاً أن حدساً بعينه، سمي من بعد «J.R.»،
يستلزم امتناع القرار في مسألة هيلبرت العاشرة
\citep{DavisPutnamRobinson1961}. وظلت روبنسون تعالج المسألة طوال ستينيات
القرن العشرين، حتى أثبت الرياضي الروسي الشاب يوري ماتيياسيفيتش سنة 1970
فرضية J.R. وأصبحت النتيجة المؤلفة تعرف بمبرهنة
ماتيياسيفيتش--روبنسون--ديفيس--بوتنام، أو اختصاراً مبرهنة MRDP. ثم انعقدت
بين ماتيياسيفيتش وروبنسون صداقة، واشتركا في بحوث عدة. وكتبت إليه مرة:
«يسعدني حقاً أننا، إذ نعمل معاً (وبيننا آلاف الأميال)، نحرز بوضوح تقدماً
أكبر مما كان يمكن لأي منا إحرازه منفرداً» \citep[45]{Matijasevich1992}.

وكانت روبنسون أول امرأة تتولى رئاسة الجمعية الرياضية الأمريكية، وأول
امرأة تُنتخب في الأكاديمية الوطنية للعلوم (National Academy of
Sciences). وتوفيت في 30 يوليو 1985، ولها 65 عاماً، بعدما شُخّص فيها سرطان
الدم.

\begin{reading}
وجُمعت بحوث روبنسون الرياضية في \textit{أعمالها المجموعة}
\citep{Robinson1996}، وفيها أيضاً إعادة طبع لترجمتها المنشورة في
الأكاديمية الوطنية للعلوم \citep{Feferman1994}. ونشرت أختها الكبرى
كونستانس ريد «سيرة جوليا الذاتية» مستندة إلى أحاديث أجرتها معها
\citep{Reid1986}، ثم نشرت ترجمة مستوفاة لحياتها \citep{Reid1996}. وصنع
جورج تشيشتشيري فيلماً وثائقياً قصيراً في روبنسون ومسألة هيلبرت العاشرة
\citep{Csicsery2016}. ولخبر موجز عن مشاركة يوري ماتيياسيفيتش لروبنسون
وأثرها في عمله، يُنظر \citep{Matijasevich1992}.
\end{reading}

\end{document}
